%%%%%%%%%%%%%%%%%%%%%%%%%%%%%%%%%%%%%%%%%%%%%%%%%%%%%%%%%%%%%%%%%%%%%%%%%%%%%%%
% Supplementary Material for main29.tex
%%%%%%%%%%%%%%%%%%%%%%%%%%%%%%%%%%%%%%%%%%%%%%%%%%%%%%%%%%%%%%%%%%%%%%%%%%%%%%%

\documentclass[]{seismica}
\usepackage{tabularx}
\usepackage{placeins}
\usepackage{longtable}
\usepackage{pdflscape}
\usepackage{booktabs}
\usepackage{etoolbox}

\graphicspath{
  {../data/outputs/170_generate_publication_figures/}
  {../data/outputs/145_analyze_regional_detectability/}
  {../data/outputs/140_search_for_direct_seismic_waves/}
  {../data/outputs/130_analyze_seismic_polarization/}
  {../data/outputs/120_estimate_acoustic_energetics/}
  {../data/outputs/110_measure_named_events/}
  {../data/outputs/100_measure_event_amplitudes_and_acoustic_seismic_coupling/figures/}
  {../data/outputs/100_measure_event_amplitudes_and_acoustic_seismic_coupling/}
  {../data/outputs/090_analyze_finite_distance_propagation/}
  {../data/outputs/080_analyze_planar_array_propagation/}
  {../data/outputs/060_validate_infrasound_baseline_correction/}
  {../data/outputs/040_validate_event_catalogue_with_array_processing/figures/}
  {../data/outputs/040_validate_event_catalogue_with_array_processing/}
  {../data/outputs/030_reconcile_manual_event_catalogues/}
  {../data/outputs/020_analyze_ksc_weather/}
  {figures/}
}

\newcommand{\pendingfigure}[1]{%
  \fbox{\parbox[c][50mm][c]{0.92\textwidth}{\centering\textbf{Pending regenerated product}\par#1}}}

\renewcommand{\thefigure}{S\arabic{figure}}
\renewcommand{\thetable}{S\arabic{table}}
\renewcommand{\theequation}{S\arabic{equation}}

\title{Supplementary Material for ``Close-Range Seismo-Acoustic Observations of the 2016 Falcon 9 AMOS-6 Failure at Cape Canaveral''}
\shorttitle{Supplementary material: 2016 Falcon 9 failure}

\author[1]{Glenn Thompson\thanks{Corresponding author: \texttt{thompsong@usf.edu}; ORCID: 0000-0002-9173-0097}}
\author[2,3]{Robert G. Brown}
\author[1,4]{Alexandra Farrell}
\author[2,3]{John Kiriazes}
\author[1]{Stephen R. McNutt}
\author[1,5]{Christopher Mehta}
\affil[1]{School of Geosciences, University of South Florida, Tampa, Florida, USA}
\affil[2]{NASA Kennedy Space Center, Florida, USA}
\affil[3]{Present address: Florida Space Institute, University of Central Florida, Orlando, Florida, USA}
\affil[4]{Present address: University of Alaska Fairbanks, Fairbanks, Alaska, USA}
\affil[5]{Present address: Howard University, Washington, DC, USA}

\begin{document}

\makeseistitle{
  \begin{summary}{Supplementary overview}
  This Supplementary Material preserves detailed processing definitions,
  sensitivity tests, regional screening results, named-event waveforms,
  finite-distance analysis, polarization overviews, and acoustic-energetics
  products supporting the shortened main manuscript. The complete event
  catalogue is tabulated below, and machine-readable event, station, and
  measurement tables will accompany the archived workflow.
  \end{summary}
}

\section{Supplementary methods}

\subsection{S1. Response correction and baseline validation}
\label{sec:supp_response}

The three seismic components were corrected using the manufacturer response
for the Nanometrics Trillium Compact Posthole seismometer and Centaur digitizer.
The combined sensitivity was $3.0172\times10^8$ counts
$(\mathrm{m\,s^{-1}})^{-1}$. With the Centaur operated at a 40-V peak-to-peak
input range, the digitizer sensitivity was $4.0\times10^5$ counts~V$^{-1}$ and
the implied seismometer sensitivity approximately
754~V~$(\mathrm{m\,s^{-1}})^{-1}$.

The pressure channels were converted independently using empirical
sensor-specific field factors. These were derived from repeated approximately
sinusoidal 1-m vertical-displacement tests, expected to produce approximately
12~Pa peak-to-peak at sea level, and door-slam tests comparing sensitivity and
waveform fidelity among 25 USF infraBSU sensors. Approximately $\pm20\%$ is
the estimated field-calibration uncertainty.
% Verify the notably lower DD2 counts-per-Pa factor against the archived
% calibration record before submission. If valid, document the sensor/gain
% difference in the repository calibration note.

Before response correction, the median count level from 20 to 2~s before the
opening signal was removed from each pressure channel. All traces received a
5\% cosine taper. Frequency-domain response removal used a 0.02, 0.05, 80, and
100~Hz prefilter and a 60-dB water level. No additional Butterworth high-pass
filter was applied to records used for physical pressure measurements.

Slow pressure variation was removed using a centred 1.0-s moving median,
containing 251 samples at 250~Hz:
\begin{equation}
 p_{i,\mathrm{corr}}(t)=p_i(t)-b_i(t).
\end{equation}
The calculation is symmetric and noncausal but does not impose a systematic
causal delay. Independent recomputation reproduced the canonical processed
stream sample-for-sample, while all seismic samples remained unchanged.

Sensitivity was tested with 0.125, 0.25, 0.5, 1.0, and 2.0-s median windows on
24 clear three-channel events. Median-stack peak-to-peak pressures were 107.6
and 107.8~Pa for 0.25 and 0.5~s, compared with 90.2 and 90.6~Pa for 1.0 and
2.0~s. The 1--2-s values form a stable long-window plateau but are
approximately 16\% below those from 0.25--0.5~s. The 0.125-s window shortened
the inferred median duration, indicating partial absorption of the impulse.
Consequently, all reported physical amplitudes use the same 1.0-s correction,
and window duration remains a processing uncertainty.

An experimental automatic compression--rarefaction algorithm succeeded for
92 of 96 waveforms from the 24 events. All four failures belonged to the
largest complex event, showing that a simple positive-then-negative pulse model
is not sufficiently general to replace manual review of key events.

\FloatBarrier
\subsection{S2. Geometry, weather, and timing coordinates}
\label{sec:supp_weather}

Geographic source--receiver distances were calculated on WGS84 and array
coordinates transformed to UTM Zone~17N. The source-to-seismic-site distance
was 1419.9~m and its SLC-40 forward azimuth 19.44$^\circ$, corresponding to a
199.44$^\circ$ back azimuth. DD1, DD2, and DD3 distances were 1439.96,
1408.34, and 1412.85~m. On 6 October 2016, as Hurricane Matthew approached
KSC, co-author Robert G. Brown removed the three above-ground pressure sensors,
electronics enclosures, and solar panels as a precaution; the buried
seismometer remained in place. The pressure sensors had not been formally
surveyed before removal, so their locations were reconstructed from cable ends
surveyed subsequently. Consequently, 1--2-m relative errors are more credible
than the formal centimetre precision of the survey.

The meteorological analysis used 105 records from three KSC towers at seven
5-min epochs between 13:05 and 13:35~UTC. Measurement elevations spanned
16.5--139.9~m. Wind observations were converted from meteorological direction
to east and north vectors, averaged within 10-m elevation bins, and projected
onto the SLC-40--BCHH direction. The nominal 0--65-m source-height summary is
an estimate because the lowest valid wind observations were near 40~m.

Atmospheric conditions changed little during the 30-min interval
(Figures~\ref{fig:supp_weather_wind} and
\ref{fig:supp_weather_time}). Mean air temperature was
$26.4\pm0.3\,^{\circ}$C and mean relative humidity was
$83.9\pm3.7\%$, where the uncertainties are sample standard deviations over
all available observations. Across the complete wind dataset, the vector-mean
wind was 5.4~m~s$^{-1}$ from 144$^\circ$. Vector-mean speeds at the seven
observation epochs ranged only from 5.3 to 5.5~m~s$^{-1}$, with a temporal
standard deviation of 0.1~m~s$^{-1}$; vector-mean directions ranged from
142$^\circ$ to 149$^\circ$.

The still-air speed was
\begin{equation}
 c_0=331.3+0.606T+1.26h,
\end{equation}
where $T$ is $^\circ$C and $h$ relative humidity expressed as a fraction. Mean
conditions gave $c_0=348.3$~m~s$^{-1}$. The estimated lower-profile
vector-mean wind was 4.6~m~s$^{-1}$ from 145$^\circ$; its path-parallel
component was +2.6~m~s$^{-1}$, giving
$c_{\mathrm{eff}}=351.0$~m~s$^{-1}$. The corresponding travel time over
1419.9~m was 4.046~s, compared with 4.076~s in still air. Thus, the wind
correction shortened the predicted travel time by approximately 0.031~s.

\begin{figure*}[!tbp]
  \centering
  \includegraphics[width=0.96\textwidth]{supplement_weather_wind_structure.pdf}
  \caption{Meteorological wind structure during 13:05--13:35~UTC.
  (a) Distribution of valid tower observations by meteorological wind
  direction and speed. (b) Individual wind-speed observations and vector
  means within 10-m height bins. (c) Corresponding wind directions.
  (d) Height-dependent wind component projected onto the SLC-40--BCHH
  propagation direction; positive values indicate wind toward BCHH. Dashed
  horizontal lines mark the nominal 65-m upper limit used for the
  source-height wind summary. Because valid wind measurements began near
  40~m, the lowest part of the profile is extrapolated from the lowest
  available observations.}
  \label{fig:supp_weather_wind}
\end{figure*}

\begin{figure*}[!tbp]
  \centering
  \includegraphics[width=0.96\textwidth]{supplement_weather_temporal_variability.pdf}
  \caption{Atmospheric conditions at seven 5-min epochs from 13:05 to
  13:35~UTC. (a) Spatially averaged vector-mean and scalar-mean wind speeds.
  (b) Vector-mean meteorological wind direction. (c) Mean air temperature.
  (d) Mean relative humidity. Shading in panels (c) and (d) denotes one
  spatial standard deviation among the available observations at each epoch.
  The small changes demonstrate that the propagation correction was stable
  throughout the impulsive-event sequence.}
  \label{fig:supp_weather_time}
\end{figure*}

\FloatBarrier

Observed UTC, DD2-referenced reduced arrival time, and source time were kept
separate. Differential reductions for DD1, DD2, and DD3 were 0.090107, 0, and
0.012867~s. They remove only array-aperture travel-time differences. Named
source times subsequently subtract the full source-to-DD2 travel time. The
rounded 351~m~s$^{-1}$ meteorological value is used for source-to-array
reduction; the independently observed approximately 352~m~s$^{-1}$ array
velocity is a result, not an input to the chronology.

\FloatBarrier
\subsection{S3. Catalogue association and waveform review}
\label{sec:supp_catalogue}

Reconciliation began with 890 manual picks, of which 384 were accepted
impulsive pressure arrivals. Channel names were mapped to DD1--DD3 and reduced
to DD2. Coincidence windows from 0.008 to 0.060~s were tested. The number of
multichannel candidates stabilized at 154 from 0.040 to 0.060~s; 0.040~s, the
smallest value on the plateau, was adopted. An event required at least two
pressure channels.

One marginal DD2--DD3 event was retained because correlation was 0.879 and
robust peak SNR 8.0. One DD1--DD2 candidate was rejected despite SNR 5.19
because no coherent DD3 signal was visible and targeted DD2--DD3 correlation
was only 0.092. The final catalogue contains 153 accepted transients and 369
accepted pressure picks: 90 events on DD2--DD3 and 63 on all three channels.
No accepted event was represented only by DD1--DD2.

Paired unshifted and differentially aligned displays were generated for every
retained event. Corrected pick spans had a median of 4.87~ms, interquartile
range 2.13--8.87~ms, and maximum 36.24~ms. These values provide an internal
check on event association and differential travel-time correction. Aligned
review stacks were diagnostics only; later standardized physical measurements
remain authoritative.

Standardized quality measurements were available for all 153 events. Robust
peak SNR had a median of 7.86 and interquartile range 5.49--16.10; median RMS
SNR was 2.23. Median DD2--DD3 correlation was 0.876, with first quartile
0.718, and median residual DD2--DD3 lag after geometric correction was zero.
Aligned review stacks had median absolute peaks of 16.7~Pa for the mean stack
and 15.4~Pa for the median stack. Their largest peak was approximately
1.386~kPa. Because those values came from short visual-review windows, they
were not substituted for the later sensor-specific amplitude measurements.

\begin{figure*}[!tbp]
  \centering
  \includegraphics[width=0.88\textwidth]{reduced_time_event_count_vs_window.pdf}
  \caption{Sensitivity of reduced-time association to the coincidence window.
  The number of multichannel candidates stabilized at 154 for windows from
  0.040 to 0.060~s. The adopted 0.040-s window is the smallest value on this
  plateau. Subsequent waveform review rejected one ambiguous candidate,
  producing the final catalogue of 153 accepted transients.}
  \label{fig:supp_association_window}
\end{figure*}

\begin{table*}[!tbp]
\centering
\caption{Nested event-quality definitions. Catalogue membership is preserved
independently of these waveform and array classes.}
\label{tab:supp_quality}
\begin{tabularx}{\textwidth}{lXr}
\hline
Subset & Definition and intended use & Count \\
\hline
Catalogue & Accepted multichannel transient; chronology and event rate & 153 \\
Usable waveform & Robust peak SNR $\geq3$ and DD2--DD3 correlation $\geq0.5$ & 138 \\
Core waveform & Robust peak SNR $\geq5$ and DD2--DD3 correlation $\geq0.7$ & 107 \\
High-quality planar & Coherence $\geq0.90$, window-stability
$\sigma_{\mathrm{BAZ}}\leq0.5^\circ$, and
$\sigma_v\leq3.0$~m~s$^{-1}$ & 51 \\
Intermediate planar & Remaining solutions with coherence $\geq0.75$,
$\sigma_{\mathrm{BAZ}}\leq1.0^\circ$, and
$\sigma_v\leq8.0$~m~s$^{-1}$ & 45 \\
Weak or broad planar & Remaining time-domain solutions & 57 \\
Regression & Independent pressure and seismic SNR plus complete amplitude capture & 46 \\
\hline
\end{tabularx}
\end{table*}

\FloatBarrier
\subsection{S4. Array processing and finite-range sensitivity}
\label{sec:supp_array}

The catalogue was evaluated independently with frequency-domain Bartlett
beamforming in InfraPy \citep{webster_infrapy_2025}. The primary event-centred analysis used 0.5-s windows,
a 2--20-Hz band, a 0.5$^\circ$ back-azimuth grid, and a 2~m~s$^{-1}$
trace-velocity grid. The empirical coherence threshold was the 99.5th
percentile of 153 pre-event background windows, corresponding to
$F=3.02$. Search-grid-edge solutions were flagged rather than silently
discarded.

In the event-centred analysis, 145 of 153 catalogue entries (94.8\%) were
within 5$^\circ$ of SLC-40 and 151 (98.7\%) were within 10$^\circ$. Median
back azimuth, apparent trace velocity, and $F$ were 200.67$^\circ$,
354.6~m~s$^{-1}$, and 5.04. In total, 113 entries exceeded the empirical
background threshold, 77 had $F\geq5$, and 42 had $F\geq10$. Two solutions
reached a back-azimuth grid edge and seven reached a trace-velocity edge.
Apparent-velocity scatter decreased from 30.7~m~s$^{-1}$ for the complete
catalogue to 12.1~m~s$^{-1}$ for entries with $F\geq10$, showing that the
largest deviations were concentrated among weakly coherent windows.

\begin{figure*}[!tbp]
  \centering
  \includegraphics[width=0.96\textwidth]{event_centred_bartlett_validation.pdf}
  \caption{Event-centred InfraPy Bartlett validation. (a) Distributions of the
  $F$ statistic for catalogue-centred and independent pre-event windows; the
  dashed line is the 99.5th-percentile background threshold. (b) Back azimuth
  and apparent trace velocity, coloured by $F$; dashed and dotted reference
  lines show the geometric SLC-40 direction and the 351.0~m~s$^{-1}$
  meteorological prediction. Red crosses identify search-grid-edge solutions.
  (c) $F$ versus fixed-geometry robust stack peak SNR. (d) $F$ versus
  DD2--DD3 waveform correlation. Colours in (c) and (d) show absolute
  back-azimuth error, clipped at 10$^\circ$.}
  \label{fig:supp_bartlett_event}
\end{figure*}

The blind continuous scan used the same 0.5-s windows and 2--20-Hz band at
0.1-s steps, but a coarser 1$^\circ$ and 5~m~s$^{-1}$ screening grid. It
analysed 17,839 overlapping windows without using catalogue times to position
them. InfraPy grouped the above-threshold windows into 12 broader activity
episodes, while an independent local-maximum search found 418 blind $F$ peaks.
A greedy one-to-one comparison associated 129 of the 153 catalogue entries
(84.3\%) with a peak within 0.25~s. The remaining 24 catalogue entries were
generally weak, closely spaced, or insufficiently coherent to form a distinct
maximum under the fixed analysis settings.

Conversely, 289 local maxima were unmatched. They are not an independent count
of additional explosions: overlapping windows, ringing, and extended signals
can generate several maxima from one physical source. Of these, 184 passed
broad direction, trace-velocity, and grid-edge screens. That large residual
demonstrates that these permissive tests are insufficient to establish
additional events; confirmation would require waveform review or independent
network evidence.

\begin{figure*}[!tbp]
  \centering
  \includegraphics[width=0.97\textwidth]{continuous_full_beam_results.pdf}
  \caption{Blind continuous InfraPy Bartlett scan. (a) $F$ statistic through
  time, with the empirical background threshold, local maxima, and the 12
  broader activity episodes. (b) Back azimuth and (c) apparent trace velocity
  for above-threshold windows only, coloured by $F$. Low-$F$ grid maxima are
  suppressed because their direction and velocity are not stable physical
  measurements; red crosses mark above-threshold grid-edge solutions. The
  local maxima and activity episodes describe coherent array activity rather
  than counts of physical explosions.}
  \label{fig:supp_bartlett_continuous}
\end{figure*}

Blind-scan recovery increased from 129 of 153 entries (84.3\%) for the full
catalogue to 75 of 77 (97.4\%) for entries with event-centred $F\geq5$ and
41 of 42 (97.6\%) for $F\geq10$. This confirms that non-recovery was
concentrated among weak or overlapping signals. Notebook~040 also applied the
same nominal SNR and correlation thresholds used for the usable and core
definitions to an independently recomputed fixed-window metric set. Those
calculations selected 132 and 100 entries, rather than the authoritative 138
and 107 in Table~\ref{tab:supp_quality}, because their waveform windows and
stack construction differed. They are retained as processing diagnostics but
are not treated as replacements for the catalogue quality classes.

Parameter sensitivity used 0.25-s windows at 4--30~Hz, the primary 0.5-s
windows at 2--20~Hz, and 1.0-s windows at 1--10~Hz. The primary configuration
gave the smallest median absolute back-azimuth error (1.43$^\circ$), the
largest median $F$, and only seven velocity-grid-edge solutions. The 0.25-s
configuration gave a median apparent velocity of 369.9~m~s$^{-1}$ and 33
velocity-edge solutions; the 1.0-s configuration gave 354.7~m~s$^{-1}$ but 49
velocity-edge solutions and a median absolute azimuth error of 2.82$^\circ$.
Because the distribution of $F$ changes with window and passband, absolute
$F$ values were interpreted relative to each configuration's background
rather than compared without qualification.

\begin{figure*}[!tbp]
  \centering
  \IfFileExists{../data/outputs/040_validate_event_catalogue_with_array_processing/figures/parameter_sensitivity_summary.pdf}{%
    \includegraphics[width=0.94\textwidth]{parameter_sensitivity_summary.pdf}%
  }{%
    \pendingfigure{Bartlett window-and-band parameter sensitivity. Re-run
    Notebook 040 to generate this new product.}%
  }
  \caption{Sensitivity of event-centred Bartlett results to window length and
  frequency band. (a) Median absolute back-azimuth error. (b) Median apparent
  trace velocity; the dotted line is the meteorological prediction.
  (c) Numbers of back-azimuth and trace-velocity search-grid-edge solutions.
  The 0.5-s, 2--20-Hz configuration provides the best overall compromise.}
  \label{fig:supp_bartlett_sensitivity}
\end{figure*}

\FloatBarrier

The time-domain inversion classified solutions independently of waveform
morphology. High-quality solutions required maximum coherence $\geq0.90$,
$\sigma_{\mathrm{BAZ}}\leq0.5^\circ$, and
$\sigma_v\leq3.0$~m~s$^{-1}$ across nine scoring-window perturbations;
intermediate solutions used corresponding thresholds of 0.75, 1.0$^\circ$,
and 8.0~m~s$^{-1}$. This yielded 51 high-quality, 45 intermediate, and 57
weak or broad solutions. High-quality events gave a
coherence-weighted back azimuth of 200.73$^\circ$ with 0.70$^\circ$
dispersion and apparent velocity of 352.0~m~s$^{-1}$ with
3.85~m~s$^{-1}$ dispersion. The principal explosion yielded back azimuth
200.3$^\circ$, velocity 352.5~m~s$^{-1}$, and coherence 0.9984. Its scoring-
window sensitivity was 0.082$^\circ$ and 0.41~m~s$^{-1}$; these values
measure repeatability, not formal spatial resolution.

Time-domain planar processing aligned interpolated pressure waveforms over a
grid of back azimuth and apparent speed. A circular-wavefront test instead
used exact horizontal distances from trial source points. With receiver
coordinates $\mathbf{x}_i$, source $\mathbf{s}$, and speed $v$,
\begin{equation}
 \Delta t_i=\frac{\lVert\mathbf{x}_i-\mathbf{s}\rVert-
 \lVert\mathbf{x}_{\mathrm{DD2}}-\mathbf{s}\rVert}{v}.
\end{equation}
Receiver elevations were unavailable, so this was a horizontal 2-D model.

The median planar and fixed-SLC-40 speeds were 352.5 and
355.5~m~s$^{-1}$, respectively, while the median event-by-event speed
difference was +4.0~m~s$^{-1}$. Fixed-source coherence was lower by a median
of 0.00512. The speed change is best interpreted as compensation for the
approximately 1$^\circ$ directional offset imposed by nominal coordinates,
not faster atmospheric propagation. Sixty-one best-fitting along-range
positions reached a search-grid boundary, and all 153 along-range profiles
remained open at a coherence decrease of 0.01. The finite-distance inversion
therefore did not resolve source range. Cross-range offsets should likewise be
treated as model sensitivity rather than physical source locations.

\begin{figure*}[!tbp]
  \centering
  \includegraphics[width=0.96\textwidth]{supplement_planar_vs_fixed_source_and_range_sensitivity.pdf}
  \caption{Planar and finite-distance sensitivity. (a) Planar apparent speed
  versus fixed-SLC-40 speed. (b) Corresponding coherence. (c) Median
  event-relative coherence deficit over trial source positions. The open
  along-range ridge demonstrates unresolved source range.}
  \label{fig:supp_finite}
\end{figure*}

\FloatBarrier
\subsection{S5. Amplitude, spectra, transfer, and energetics}
\label{sec:supp_energy}

Notebook~070 generated fixed six-channel event segments extending from 0.50~s
before to 0.60~s after each DD2-referenced catalogue time. Subsequent amplitude
measurements used narrower event-specific measurement intervals within these
segments. Pressure channels were differentially aligned; seismic
channels were not shifted. Peak-to-peak pressure was the median of individual
sensor measurements, not the amplitude of a median-stacked waveform.
Three-component vector PGV was
\begin{equation}
 \mathrm{PGV}_{3C}=\max_t\sqrt{v_Z^2(t)+v_N^2(t)+v_E^2(t)}.
\end{equation}
Events entered regression only if pressure and seismic SNR thresholds were
met and the peak was not clipped by a window boundary. Ordinary least squares,
equal-weight orthogonal regression, leave-one-out fits, exclusion of the
principal explosion, and event bootstrap assessed robustness.

Event spectra and cross-spectra used pressure aligned to DD2 and vertical
ground velocity. The H1 estimator, coherence, and bootstrap intervals were
supplemented by mismatched-event and time-shift permutation tests. The
continuous Phase~I estimate used nonoverlapping 10-s segments. Removing the
largest events, equalizing pressure RMS, and varying the pressure-beam
construction preserved the 4--5 and 20--24-Hz bands.

For named events, a short pre-signal interval removed residual constant offset
and positive, negative, and peak-to-peak pressures were measured independently
on DD1--DD3. The array summary is the median of scalar sensor measurements.
Pressures reduced to 1~km use $p_{1\,\mathrm{km}}=p_{\mathrm{obs}}r/1000$ and
are diagnostic geometric reductions, not direct measurements at 1~km.

Sound exposure level is
\begin{equation}
 \mathrm{SEL}=10\log_{10}\left[
 \frac{\int p^2(t)\,dt}{p_0^2(1~\mathrm{s})}\right],
 \qquad p_0=20~\mu\mathrm{Pa}.
\end{equation}
Peak pressure level, used separately from SEL, is
\begin{equation}
 L_{\mathrm{peak}}=20\log_{10}\left(\frac{|p_+|}{20~\mu\mathrm{Pa}}\right).
\end{equation}
Acoustic energy assumes hemispherical radiation without attenuation. It is
converted to a radiated acoustic TNT equivalent by division by
$4.184\times10^6$~J~kg$^{-1}$. Illustrative total-yield scenarios additionally
divide by acoustic efficiency $\eta_{\mathrm{ac}}$.

The upper-stage event had median positive, negative, and peak-to-peak
pressures of +38.8, $-25.4$, and 77.1~Pa; its corresponding 1-km reductions
were +54.9, $-35.8$, and 108.6~Pa. The principal explosion had medians of
+1.392~kPa, $-157$~Pa, and 1.549~kPa. Across DD1--DD3 its positive pressure
ranged from 1.31 to 1.51~kPa and peak-to-peak pressure from 1.44 to 1.67~kPa;
1-km median reductions were 1.967 and 2.189~kPa for positive and peak-to-peak
pressure. Payload impact had medians of +279.1, $-115.3$, and 390.7~Pa, while
the post-impact explosion had +297.0, $-113.4$, and 410.6~Pa.

Across the full catalogue, median stack peak-to-peak pressure was 29.5~Pa
(interquartile range 16.2--51.3~Pa) and median vector PGV
$5.6\times10^{-5}$~m~s$^{-1}$. The full-catalogue median dimensional
pressure-to-vector-PGV ratio was
$4.58\times10^5$~Pa~m$^{-1}$~s, compared with
$4.34\times10^5$~Pa~m$^{-1}$~s in the 46-event regression subset. Scatter
increased sharply below vector PGV of approximately $10^{-4}$~m~s$^{-1}$,
coincident with decreasing seismic SNR and incomplete peak capture.

The 46-event log--log regression gave slope 0.962 with a 95\% event-bootstrap
interval of 0.893--1.009, $R^2=0.953$, and Spearman correlation 0.931.
Equal-weight orthogonal regression gave 0.986; excluding the principal
explosion gave 0.958 with $R^2=0.937$; and leave-one-out slopes ranged from
0.952 to 0.971. The near-unit scaling was therefore not imposed by the
principal explosion.

Event coherence was approximately 0.61--0.68 near 4.5--5.3~Hz and
0.43--0.70 near 21--28~Hz, above frequency-specific permutation thresholds.
The continuous Phase~I H1 estimate reached approximately
$1.1\times10^{-5}$~(m~s$^{-1}$)~Pa$^{-1}$ near 4.7~Hz and
$5$--$7\times10^{-6}$~(m~s$^{-1}$)~Pa$^{-1}$ near 21--25~Hz. These broad
features persisted after removing the strongest events, equalizing pressure
RMS, and changing pressure-beam construction.

\begin{figure*}[!tbp]
  \centering
  \IfFileExists{../data/outputs/100_measure_event_amplitudes_and_acoustic_seismic_coupling/figures/supplement_transfer_robustness_tests.pdf}{%
    \includegraphics[width=0.94\textwidth]{supplement_transfer_robustness_tests.pdf}%
  }{%
    \pendingfigure{Transfer-function robustness tests from Notebook 100,
    including exclusion of dominant events, equal event weighting, and
    alternative pressure-beam definitions.}%
  }
  \caption{Sensitivity of the empirical pressure-to-vertical-velocity transfer
  estimate. The tests exclude progressively more of the strongest events and
  vary event weighting and pressure-beam construction. Persistence of the
  broad 4--5- and 20--24-Hz features indicates that they are not produced
  solely by the principal explosion or by one stacking convention.}
  \label{fig:supp_transfer_robustness}
\end{figure*}

Selected event windows were converted to absolute DD2 intervals and all
overlaps merged. The continuous pressure stream was aligned and $p^2(t)$
integrated once per disjoint event complex, separately by channel. Cumulative
energy is the channel-wise sum of these nonoverlapping complex energies. The
selection covers the union of high-quality and named-event intervals, not
necessarily every contribution during the full 26-min sequence.

\begin{table*}[!tbp]
\centering
\small
\caption{Acoustic energetics of the four manually reviewed opening-sequence
windows. SEL is referenced to 20~$\mu$Pa and 1~s. $E_{2\pi}$ assumes
hemispherical radiation; the final column is an illustrative 1\%-efficiency
scenario, not a calibrated blast yield.}
\label{tab:supp_energetics}
\begin{tabular}{lrrrr}
\hline
Signal & SEL (dB) & $E_{2\pi}$ (J) & Acoustic TNT (kg) &
Yield at 1\% (t TNT) \\
\hline
Upper-stage event       & 114.21 & $3.21\times10^6$ & 0.768 & 0.0768 \\
Principal explosion     & 142.54 & $2.20\times10^9$ & 525   & 52.5 \\
Payload impact          & 126.44 & $5.40\times10^7$ & 12.9  & 1.29 \\
Post-impact explosion   & 127.05 & $6.22\times10^7$ & 14.9  & 1.49 \\
\hline
\end{tabular}
\end{table*}

Table~\ref{tab:supp_energetics} uses the manually reviewed named-event windows
defined in Notebook~110 and integrated in Notebook~120. The independent
event-complex calculation in Notebook~100 instead merges overlapping catalogue
windows before integration. Under that definition, the lower-stage complex
contained a median $1.97\times10^9$~J across DD1--DD3 (range
$1.68$--$2.25\times10^9$~J), equivalent to 471~kg TNT as radiated acoustic
energy. The difference from the $2.20\times10^9$-J named-window value reflects
the different integration and local-baseline definitions; it is not an
independent estimate of a second physical explosion.

\begin{figure*}[!tbp]
  \centering
  \IfFileExists{../data/outputs/120_estimate_acoustic_energetics/key_event_pressure_and_cumulative_acoustic_energy.pdf}{%
    \includegraphics[width=0.96\textwidth]{key_event_pressure_and_cumulative_acoustic_energy.pdf}%
  }{%
    \pendingfigure{Channel-specific named-event pressure waveforms and
    cumulative hemispherical acoustic-energy curves from Notebook 120.}%
  }
  \caption{Channel-specific pressure and cumulative hemispherical acoustic
  energy for the four manually reviewed opening-sequence windows. Time is
  relative to the predicted arrival at each sensor, and grey shading marks the
  integration interval. The cumulative curves provide a direct check on
  whether each selected window captures the principal energy contribution.}
  \label{fig:supp_named_energy_windows}
\end{figure*}

The 51 selected high-quality or named catalogue windows merged into 47
nonoverlapping event complexes. Summing those complexes independently on each
channel gave a median cumulative hemispherical acoustic energy of
$2.27\times10^9$~J (range $1.95$--$2.55\times10^9$~J), equivalent to 543~kg
TNT as radiated acoustic energy (range 467--610~kg). The corresponding median
illustrative total-energy equivalent is 54.3~t TNT at 1\% acoustic efficiency;
varying efficiency from 10\% to 0.1\% gives 5.43--543~t. These quantities cover
only the union of the selected intervals and are not estimates of all acoustic
energy radiated during the complete sequence.

\begin{figure*}[!tbp]
  \centering
  \IfFileExists{../data/outputs/100_measure_event_amplitudes_and_acoustic_seismic_coupling/figures/deduplicated_event_complex_and_cumulative_energy.pdf}{%
    \includegraphics[width=0.94\textwidth]{deduplicated_event_complex_and_cumulative_energy.pdf}%
  }{%
    \pendingfigure{De-duplicated event-complex energies and running cumulative
    acoustic energy. Regenerate Notebook 100 before finalizing the cumulative
    numerical result.}%
  }
  \caption{Selected nonoverlapping event-complex acoustic energies and their
  cumulative sum. Overlapping catalogue windows are integrated only once.
  The cumulative value refers only to the union of the selected measurement
  intervals, rather than to all acoustic energy radiated during the 26-min
  sequence.}
  \label{fig:supp_cumulative}
\end{figure*}

\subsubsection*{Complete event measurements}

Tables~\ref{tab:supp_event_propagation} and
\ref{tab:supp_event_amplitudes} report a compact set of measurements for every
accepted catalogue entry. They are read directly from
\texttt{event\_acoustic\_seismic\_amplitudes.csv}, so the printed tables and
machine-readable product have the same provenance. Time is measured from the
first accepted DD2-referenced arrival. The pressure values in these tables are
fixed-window aligned-stack measurements intended for comparisons across the
full catalogue; the separately reviewed scalar sensor medians in
Table~\ref{tab:supp_energetics} remain authoritative for the four named events.

\IfFileExists{../data/outputs/100_measure_event_amplitudes_and_acoustic_seismic_coupling/event_acoustic_seismic_amplitudes.csv}{%
\DTLloaddb{eventmetrics}{../data/outputs/100_measure_event_amplitudes_and_acoustic_seismic_coupling/event_acoustic_seismic_amplitudes.csv}

\begin{landscape}
\scriptsize
\setlength{\tabcolsep}{4pt}
\begin{longtable}{rrrrlrrr}
\caption{Chronology, association, and time-domain planar-wave measurements for
all 153 accepted pressure transients. Time is elapsed from the first
DD2-referenced catalogue arrival. Span is the corrected difference between the
first and last associated pressure picks. Classes A, B, and C denote
high-quality, intermediate, and weak or broad planar solutions,
respectively.}\label{tab:supp_event_propagation}\\
\toprule
Event & Time (s) & Channels & Span (s) & Class & BAZ ($^\circ$) &
$c_{\mathrm{app}}$ (m s$^{-1}$) & Coherence \\
\midrule
\endfirsthead
\multicolumn{8}{c}{Table~\thetable\ continued}\\
\toprule
Event & Time (s) & Channels & Span (s) & Class & BAZ ($^\circ$) &
$c_{\mathrm{app}}$ (m s$^{-1}$) & Coherence \\
\midrule
\endhead
\midrule
\multicolumn{8}{r}{Continued on next page}\\
\endfoot
\bottomrule
\endlastfoot
\DTLforeach*{eventmetrics}{%
  \EventNumber=event_number,%
  \ElapsedTime=elapsed_time_s,%
  \ChannelCount=channel_count,%
  \PickSpan=corrected_pick_span_s,%
  \QualityClass=solution_quality_class,%
  \BackAzimuth=best_back_azimuth_deg,%
  \ApparentSpeed=best_apparent_speed_mps,%
  \ArrayCoherence=maximum_coherence_score}{%
  \EventNumber &
  \num[round-mode=places,round-precision=2]{\ElapsedTime} &
  \ChannelCount &
  \num[round-mode=places,round-precision=4]{\PickSpan} &
  \ifdefstring{\QualityClass}{A_high_quality}{A}{%
    \ifdefstring{\QualityClass}{B_probable}{B}{C}} &
  \num[round-mode=places,round-precision=1]{\BackAzimuth} &
  \num[round-mode=places,round-precision=1]{\ApparentSpeed} &
  \num[round-mode=places,round-precision=3]{\ArrayCoherence} \\
}
\end{longtable}
\end{landscape}

\begin{landscape}
\scriptsize
\setlength{\tabcolsep}{4pt}
\begin{longtable}{rrrrrrc}
\caption{Acoustic and seismic amplitude measurements for all 153 accepted
pressure transients. Acoustic amplitude is the aligned-stack peak-to-peak
pressure. PGV is the maximum three-component vector ground velocity. Capture
is the ratio between PGV in the adopted measurement window and that in the
complete fixed segment. The final column identifies the 46 events admitted to
the pressure--PGV regression.}\label{tab:supp_event_amplitudes}\\
\toprule
Event & $p_{\mathrm{p2p}}$ (Pa) & PGV (m s$^{-1}$) & Acoustic SNR &
Seismic SNR & Capture & Regression \\
\midrule
\endfirsthead
\multicolumn{7}{c}{Table~\thetable\ continued}\\
\toprule
Event & $p_{\mathrm{p2p}}$ (Pa) & PGV (m s$^{-1}$) & Acoustic SNR &
Seismic SNR & Capture & Regression \\
\midrule
\endhead
\midrule
\multicolumn{7}{r}{Continued on next page}\\
\endfoot
\bottomrule
\endlastfoot
\DTLforeach*{eventmetrics}{%
  \EventNumber=event_number,%
  \PeakToPeak=acoustic_stack_peak_to_peak_pa,%
  \VectorPGV=pgv_vector_mps,%
  \AcousticSNR=acoustic_stack_p2p_snr,%
  \SeismicSNR=pgv_vector_snr,%
  \CaptureFraction=pgv_window_capture_fraction,%
  \RegressionFlag=passes_regression_quality}{%
  \EventNumber &
  \num[round-mode=places,round-precision=2]{\PeakToPeak} &
  \num[exponent-mode=scientific,round-mode=figures,round-precision=3]{\VectorPGV} &
  \num[round-mode=places,round-precision=1]{\AcousticSNR} &
  \num[round-mode=places,round-precision=1]{\SeismicSNR} &
  \num[round-mode=places,round-precision=2]{\CaptureFraction} &
  \ifdefstring{\RegressionFlag}{True}{yes}{no} \\
}
\end{longtable}
\end{landscape}
}{%
\pendingfigure{The complete 153-event tables require the Notebook~100 product
\texttt{event\_acoustic\_seismic\_amplitudes.csv}. Run Notebook~100 before
compiling the final Supplementary Material.}%
}

\FloatBarrier
\subsection{S6. Regional detectability screening}
\label{sec:supp_regional}

Stations operating at the principal-explosion source time were queried from
FDSN services. Pressure channels matching BDF or similar pressure codes were
examined to 1800~km. Vertical HHZ, BHZ, and, where higher-rate data were
unavailable, LHZ channels were examined to 600~km. Cached waveform and metadata
products were used before issuing new queries. Each record was response
corrected where metadata permitted and screened within broad tropospheric,
stratospheric, thermospheric/slow-acoustic, $P$-like, and $S$-like windows.

Automatic metrics compared peak amplitude and robust RMS within each phase
window with duration-matched pre- and post-event controls separated by a guard
interval. The criteria were deliberately permissive and generated review
candidates, not detections. Manual review considered onset, duration,
background stationarity, spectral content, and consistency with other
stations. Confirmation required coherent network moveout or independent array
direction information.

The search assembled 171 station--sensor targets; 166 supplied usable
waveforms and five failed because of missing data or response information.
Automatic criteria flagged 93 phase windows at 67 stations. Most flags were
persistent narrow-band background, unrelated transients, or peaks without a
distinct onset.

At IU.DWPF, the 1-Hz pressure channel restricted analysis to 0.2--0.45~Hz. No
distinct packet occurred in the direct tropospheric window. A later maximum at
13:15:42~UTC had celerity 192.8~m~s$^{-1}$ and SNR 3.33 but lay within
persistent approximately 4-s oscillations and was rejected as background. The
vertical seismic record had no accepted $P$- or $S$-like counterpart.

Three isolated pressure packets at N4.257A, N4.S54A, and N4.R55A were the best
remaining candidates (Table~\ref{tab:supp_regional_candidates}). Their
diagnostic spherical reductions to 1~km are not source pressures because
ducting, absorption, focusing, and frequency-dependent filtering violate a
simple $1/r$ law. The 98-station northern record section showed no coherent
moveout, so none was accepted as a Falcon~9 arrival.

\begin{table*}[!tbp]
\centering
\caption{Best isolated regional pressure packets. All remain unconfirmed and
were excluded from period--yield analysis.}
\label{tab:supp_regional_candidates}
\begin{tabular}{lrrrrr}
\hline
Station/channel & Distance (km) & Azimuth ($^\circ$) & Celerity (m s$^{-1}$) & Peak (Pa) & Diagnostic $p_{1\,\mathrm{km}}$ (Pa) \\
\hline
N4.257A.EP.BDF & 380.8 & 353.6 & 192.7 & 0.352 & 134 \\
N4.S54A.EP.BDF & 1026.8 & 356.4 & 197.0 & 0.323 & 332 \\
N4.R55A.EP.BDF & 1079.0 & 2.1 & 224.7 & 0.162 & 175 \\
\hline
\end{tabular}
\end{table*}

\begin{figure*}[!tbp]
  \centering
  \includegraphics[width=0.90\textwidth]{supplement_regional_station_map_best_candidates.pdf}
  \caption{Regional pressure and vertical-seismic stations screened for a
  counterpart. Filled diamonds mark the three clearest isolated pressure
  packets; none is a confirmed association.}
  % Production note: regenerate the map to remove clipped marginal labels and
  % separate the northern candidate labels before submission.
  \label{fig:supp_regional_map}
\end{figure*}

\begin{figure*}[!tbp]
  \centering
  \includegraphics[width=0.98\textwidth]{145_northern_corridor_source_reduced_section.pdf}
  \caption{Source-reduced 0.2--1-Hz pressure section for 98 northern stations.
  Traces are independently normalized and ordered by range; the line is a
  196.4~m~s$^{-1}$ trial branch. No coherent moveout confirms the isolated
  candidates.}
  \label{fig:supp_regional_section}
\end{figure*}

\FloatBarrier
\subsection{S7. Polarization and multimedia synchronization}
\label{sec:supp_polarization}

Sequence-wide ground velocity was rotated to $Z$, radial, and transverse
components using a 199.438$^\circ$ back azimuth and filtered from 0.5 to 20~Hz.
Real covariance and analytic-signal polarization attributes were calculated in
overlapping windows. Bands were 0.5--3, 3--10, 10--25, 25--45, and
0.5--20~Hz; windows shortened from 1.0~s at low frequency to 0.2~s at high
frequency. Whole-interval medians include low-amplitude background and are
descriptive only. Median degrees of polarization were 0.522, 0.698, 0.442, 0.557, and
0.533 in the 0.5--3, 3--10, 10--25, 25--45, and 0.5--20~Hz bands,
respectively.

The focused upper-stage search used causal two-pole Butterworth filters as the
primary timing analysis and zero-phase filtering as a sensitivity test. The
search stopped at +3.45~s to limit contamination by the principal explosion at
+3.601~s. Persistent four-, five-, and six-robust-standard-deviation crossings
were compared with STA/LTA and AIC variance-change estimates. Trailing
polarization windows required amplitude above four background standard
deviations, rectilinearity $\geq0.70$, degree of polarization $\geq0.55$, and
radiality $\geq0.70$.

\begin{table*}[!tbp]
\centering
\caption{Focused upper-stage onset results. The velocity envelope includes
the causal threshold-crossing range and the zero-phase sensitivity estimate.
All apparent velocities assume that the signal originated with the
upper-stage failure.}
\label{tab:supp_onsets}
\begin{tabularx}{\textwidth}{lccccX}
\hline
Band & Causal range (s) & Causal median (s) & Zero-phase median (s) & Velocity envelope (km s$^{-1}$) & Interpretation \\
\hline
0.8--3 Hz & 1.843--1.947 & 1.863 & 1.739 & 0.73--0.82 & Robust low-frequency onset \\
3--8 Hz & 1.855--1.999 & 1.895 & 1.703 & 0.71--0.84 & Best-constrained onset \\
8--20 Hz & 2.559--2.787 & 2.659 & 2.615 & 0.51--0.56 & Later or slower packet \\
20--45 Hz & Single 4$\sigma$ pick & 2.811 & none & approximately 0.51 & Not reliable \\
\hline
\end{tabularx}
\end{table*}

\begin{figure*}[!tbp]
  \centering
  \includegraphics[width=0.96\textwidth]{multiband_polarization_comparison.pdf}
  \caption{Sequence-wide degree of polarization and ellipticity in five
  frequency bands. Black points are degree of polarization and grey points
  ellipticity. Vertical markers identify configured source times and predicted
  airwave arrivals. This overview motivates, but does not itself display, the
  focused 1.7--2.0-s causal onset measurement used in the main paper.}
  \label{fig:supp_multiband}
\end{figure*}

Amplitude-gated polarization independently confirmed organized motion in the
three lower frequency bands. The first qualifying trailing window ended at
2.095~s in the 0.8--3-Hz band, constraining the contributing signal to
1.095--2.095~s; at 1.875~s in the 3--8-Hz band, constraining it to
1.475--1.875~s; and at 2.575~s in the 8--20-Hz band, constraining it to
2.375--2.575~s. No qualifying polarization interval was recovered at
20--45~Hz. These brackets reflect the full trailing-window duration and
should not be interpreted as equally precise onset picks.

\begin{figure*}[!tbp]
  \centering
  \includegraphics[width=0.97\textwidth]
    {upper_stage_phase_discrimination.pdf}
  \caption{Amplitude-gated component-energy and \(Z\)--\(R\) phase diagnostics
  for the focused upper-stage interval. Motion after the low-frequency onset
  is concentrated principally in the radial plane. Approximately quadrature
  \(Z\)--\(R\) motion develops later in the packet, consistent with
  Rayleigh-type or mixed shallow-wave motion but insufficient for a unique
  phase assignment. Pale curves represent intervals failing the amplitude
  gate.}
  \label{fig:supp_phase_discrimination}
\end{figure*}

\begin{figure*}[!tbp]
  \centering
  \includegraphics[width=0.94\textwidth]
    {upper_stage_noise_normalized_spectrogram.pdf}
  \caption{Noise-normalized \(Z\), radial, and transverse spectrograms of the
  focused upper-stage interval using 0.4-s and 1.0-s Fourier windows. Shading
  gives amplitude in decibels above the frequency-specific pre-event median.
  Low-frequency energy increases principally on the radial and vertical
  components before the predicted ground-coupled airwave.}
  \label{fig:supp_upper_stage_spectrogram}
\end{figure*}

The USLaunchReport camera range was approximately 4.255~km. Projecting the wind
onto that path gave an effective speed of 349.8~m~s$^{-1}$ and travel time of
12.16~s. Camera audio was advanced by this travel time. BCHH pressure and
ground-velocity traces were shifted using their individual ranges and
351~m~s$^{-1}$. The ground motion in these short windows is dominated by
ground-coupled airwaves, justifying the acoustic reduction for multimodal
display; it does not imply that all seismic energy propagated acoustically.

\begin{figure*}[!tbp]
  \centering
  \includegraphics[width=0.98\textwidth]{fig03_bchh_30s_zoom_with_audio_reduced_zp_true.pdf}
  \caption{Source-aligned 30-s opening sequence with synchronized camera audio,
  DD1--DD3 pressure, and three-component ground velocity. The mapped UTC scale
  remains conditional on the upper-stage visual anchor. Audio source:
  \citet{uslaunchreport_amos6_video_2016}; used with permission.}
  \label{fig:supp_opening}
\end{figure*}

\begin{figure*}[!tbp]
  \centering
  \includegraphics[width=0.92\textwidth]{principal_explosion_waveforms.pdf}
  \caption{Calibrated pressure and ground-velocity waveforms for the principal
  explosion.}
  % Production note: regenerate plotted annotations from the same authoritative
  % named-event measurements used in the main table.
  \label{fig:supp_principal}
\end{figure*}

\begin{figure*}[!tbp]
  \centering
  \IfFileExists{../data/outputs/170_generate_publication_figures/payload_explosion_waveforms.pdf}{%
    \includegraphics[width=0.92\textwidth]{payload_explosion_waveforms.pdf}%
  }{%
    \pendingfigure{Payload-impact and post-impact-explosion waveform panels.
    Regenerate Notebook 170 so that internal labels no longer say ``capsule''.}%
  }
  \caption{Calibrated waveforms for payload impact and the post-impact
  explosion, separated by 0.568~s. The labels are based on the mapped video
  chronology.}
  \label{fig:supp_payload}
\end{figure*}

The principal explosion generated surface-wave-like ground motion that
dominated the ground-velocity record for approximately 20~s. The prolonged
signal later in Phase~I persisted from approximately 13:08:48 to
13:12:00~UTC and was strongest on the seismic channels, although a pressure
contribution on DD2 cannot be excluded. Public video shows the AMOS-6 payload
descending and reaching the ground during the two payload-associated signals;
the waveforms themselves remain dominated by ground-coupled airwaves and do
not isolate a unique direct seismic impact phase.

\FloatBarrier

\section{Additional engineering and operational context}
\label{sec:supp_engineering}

SpaceX's Accident Investigation Team operated under an FAA-approved plan with
FAA oversight and participation by NASA, the U.S. Air Force, the National
Transportation Safety Board, industry experts, and other organizations
\citep{spacex_anomaly_2017,nasa_sma_falcon9}. The investigation examined more
than 3,000 telemetry and video channels, umbilical data, imagery, recovered
hardware, and subsequent tests. SpaceX reported 93~ms between the first
anomalous telemetry and loss of the second stage.

The reported initiating failure was one of three composite overwrapped pressure
vessels in the second-stage liquid-oxygen tank. Credible paths involved
supercooled liquid oxygen, potentially including solid oxygen, trapped between
an aluminium-liner buckle or void and the carbon-fibre overwrap. Fibre breakage
or friction during pressurization could ignite the oxygen. Corrective measures
returned helium loading to a warmer, previously flight-proven configuration;
the longer-term remedy redesigned the vessels to prevent liner buckling
\citep{spacex_anomaly_2017}.

SpaceX reported that the Falcon Support Building and new liquid-oxygen farm
were unaffected, the RP-1 farm largely unaffected, and pad control systems in
relatively good condition despite extensive damage elsewhere. No personnel
were injured. Falcon~9 returned to flight from Vandenberg after 135~days, while
SLC-40 returned with CRS-13 after 470~days
\citep{spacex_iridium1_2017,spacex_crs13_2017}. Activation of Launch
Complex~39A allowed Florida operations to resume during the outage.

The reported US\$196-million AMOS-6 reimbursement and approximately
US\$50-million SLC-40 restoration-and-upgrade cost exclude the launch vehicle,
schedule disruption, revenue, and complete customer and insurance consequences
\citep{magan_amos6_2016,clark_slc40_2017}. They should not be summed as a
comprehensive accident cost.

\section{Supplementary data products}

The repository deposit should contain, at minimum:
\begin{enumerate}
  \item all 153 catalogue entries with accepted picks, waveform-quality tiers,
        array solutions, amplitudes, PGV, and provenance, including the
        machine-readable source of Tables~\ref{tab:supp_event_propagation}
        and~\ref{tab:supp_event_amplitudes};
  \item the excluded-candidate audit;
  \item channel-specific named-event measurements and windows;
  \item all nonoverlapping event-complex energetic measurements and the final
        cumulative summary;
  \item the regional station/channel manifest, metadata and waveform
        availability audit, automatic metrics, candidate plots, and manual
        decisions; and
  \item the configuration, environment, and notebook-to-output manifest needed
        to reproduce every table and figure.
\end{enumerate}

% The Seismica class selects abbrvnat_seismica_upcasetitle automatically.
\bibliography{mybibfile19}

\end{document}
